\section[Cote 79 : graphes à homothéties, commutant et étoiles sur $(\mathbf Z/n)^2$]{Cote 79 : graphes à homothéties, commutant et étoiles sur $(\mathbf Z/n)^2$\protect\verifie{Folder79}}\label{sec:79}

La cote 79, \emph{Géométrie et topologie combinatoire}, compte deux cent quatre-vingt-quinze pages. Elle part des triangles d'un module libre $M$ de rang~$2$ sur un anneau commutatif $\Lambda$ : les sommets sont les éléments de $M^*$, vecteurs qui font partie d'une base, pris au signe près, et une arête de type $\pm e$ joint $\dot x$ et $\dot y$ quand $x \wedge y = \pm e$. Pour $\Lambda = \mathbf Z/n$ on obtient les polyèdres $P_{M, i}$, un par type $i \in \Lambda^*/\pm 1$. La page~19 relie deux sommets par $z = \lambda x + \lambda y$ quand $e = \lambda\, x \wedge y$ ; la page~20 cite comme connu que deux éléments de $M^*$ ont sur un anneau semi-local un partenaire de base commun, et en tire la distance au plus~$3$, \q{atteinte pour l'icosaèdre} (\lot{79}{1}{19}, \lot{79}{1}{20}). Pages~92 à~95, elle calcule le commutant de $\mathrm{SL}(M)$ et de $\mathrm{SL}(M)/\pm 1$ par le normalisateur du stabilisateur d'un point ; elle trouve l'exception $2\Lambda = \{0, 2\} \neq \{0\}$, un élément $\sigma$ induit par la matrice de lignes $(1, 0)$ et $(2, 1)$, un produit \q{direct (pas seulement ½direct !)}, et pour $\Lambda = \mathbf Z/4$ l'antipodie de l'octaèdre ; une formule biffée y donne $\sigma(\dot x) = (x + 2y)^{\cdot}$ (\lot{79}{5}{93}, \lot{79}{5}{94}, \lot{79}{5}{95}). Pages~136 à~140, le \q{fait mirifique} : $(\mathbf Z/n)^*/\pm 1$ opère sur les structures de polygone d'un ensemble à $n$ éléments, et l'axiome To~6 demande qu'une homothétie redessine chaque étoile en polygone étoilé (\lot{79}{7}{138}, \lot{79}{7}{139}). Pages~187 et~188, elle examine si une homothétie $\lambda$ d'un 2-polyèdre régulier tord la triangulation en une autre ; la condition (32) a$''$) demande que $\lambda^{-1}s$ soit à distance au plus~$3$ de $s$, \q{et en fait, égale à 3}, sauf pour l'identité et pour l'antipodie de l'octaèdre (distance~$2$) ; la distance~$1$ est exclue par les stabilisateurs, et la distance~$2$ donnerait deux sommets de même étoile. La phrase porte un mot illisible et deux lectures incertaines (\lot{79}{10}{188}). Les pages~242 à~247 traduisent un graphe homogène muni d'un groupe d'homothéties en données de théorie des groupes $(H, U_0, \sigma, T_H)$ et énoncent l'équivalence (88), avec en marge \q{[ ] résulte des autres} ; la réciproque, page~247, porte plusieurs mots illisibles (\lot{79}{13}{242}, \lot{79}{13}{245}, \lot{79}{13}{247}). Pages~275 à~279, elle demande si un sommet est déterminé par son étoile, répond sur un anneau semi-local par la même exception, et construit l'involution $a$, de déterminant $-1$ (\lot{79}{14}{276}, \lot{79}{14}{277}, \lot{79}{14}{278}).

\noindent\emph{Conventions.} Pour un groupe $H$, un sous-groupe $U_0$ et $\sigma \in H$, le graphe de classes a pour sommets $S = H/U_0$, pour origine $s_0 = U_0$ et $t_0 = \sigma U_0$ ; les arcs sont les $(hs_0, ht_0)$, $h \in H$. On note $U_1 = \sigma U_0 \sigma^{-1}$, stabilisateur de $t_0$. Les hypothèses $\sigma^2 \in U_0$ (l'élément $\sigma$ échange $s_0$ et $t_0$) et $U_0 \cap U_1 = 1$ disent que $H$ est simplement transitif sur les arcs. Un \emph{groupe d'homothéties} est un sous-groupe $Z$ du groupe des permutations de $S$ dont les éléments commutent à $H$, envoient chaque arc $(s, t)$ sur un arc $(\lambda s, \lambda^{-1} t)$, et vérifient la fidélité (72) : si $u \in H$ fixe $s_0$ et envoie $t_0$ sur $\lambda t_0$, alors $\lambda = 1$. Pour le rang~$2$, $R$ est un anneau commutatif quelconque, $M = R^2$, $x \wedge y = x_0y_1 - x_1y_0$, et $S_M = M^*/\pm 1$, où $x \in M^*$ s'il existe $y$ avec $x \wedge y = 1$. Le groupe $\mathrm{SL}(2, R)$ opère sur $S_M$, et $c \in R^*$ par $\dot x \mapsto (cx)^{\cdot}$. On dit que $R$ est \emph{exceptionnel} si $2s \in \{0, 2\}$ pour tout $s \in R$, c'est-à-dire $2R \subset \{0, 2\}$ ; avec $2 \neq 0$, c'est $2R = \{0, 2\} \neq \{0\}$.

\begin{enonce}[Pages 242 à 247, (88) et (89)\piste{79-homothety-groups-as-subgroups}]
Soient $H$, $U_0$, $\sigma$ comme ci-dessus, avec $\sigma^2 \in U_0$ et $U_0 \cap U_1 = 1$. Appelons \emph{tore} un sous-groupe $T \subset H$ tel que (a) $\sigma t \sigma^{-1} = t^{-1}$ pour $t \in T$, (b) $T$ normalise $U_0$, (c) l'application $U_0 \times T \times U_1 \to H$, $(u_0, t, u_1) \mapsto u_0tu_1$, est injective.
\begin{enumerate}
\item Pour un groupe d'homothéties $Z$, $T_H = \{u \in H \mid \exists \lambda \in Z,\ us_0 = \lambda s_0,\ ut_0 = \lambda^{-1}t_0\}$ est un sous-groupe, et c'est un tore.
\item Pour un tore $T$, les permutations $\lambda$ de $S$ qui commutent à $H$ et vérifient $\lambda s_0 = ts_0$ pour un $t \in T$, c'est-à-dire $hU_0 \mapsto htU_0$, forment un groupe d'homothéties $Z_T$.
\item Les applications $Z \mapsto T_H$ et $T \mapsto Z_T$ sont des bijections réciproques entre groupes d'homothéties et tores.
\item Un groupe d'homothéties est commutatif, et cela sans l'hypothèse $U_0 \cap U_1 = 1$.
\end{enumerate}
\end{enonce}

\begin{proof}
Une permutation qui commute à $H$ est déterminée par l'image de $s_0$, puisque $\lambda(hs_0) = h\lambda(s_0)$. Deux éléments de $H$ qui coïncident sur $s_0$ et sur $t_0$ sont égaux, puisque $U_0 \cap U_1 = 1$. Comme $\sigma^2$ est dans $U_0$ et dans $U_1$, $\sigma^2 = 1$.

Commutativité. Soient $\lambda, \mu \in Z$ et $u, v, w \in H$ tels que $u(s_0, t_0) = (\lambda s_0, \lambda^{-1}t_0)$, $v(s_0, t_0) = (\mu s_0, \mu^{-1}t_0)$, $w(s_0, t_0) = (\mu\lambda s_0, (\mu\lambda)^{-1}t_0)$ ; ils existent par la condition d'homothétie appliquée à l'arc $(s_0, t_0)$. L'élément $w^{-1}uv$ fixe $s_0$ et envoie $t_0$ sur $c\,t_0$, avec $c = \mu^{-1}\lambda^{-1}\mu\lambda$. La fidélité donne $c = 1$.

Point~1. Si $u$ correspond à $\lambda$ et $v$ à $\mu$, $uv$ correspond à $\mu\lambda$, grâce à la commutativité ; $u^{-1}$ correspond à $\lambda^{-1}$. Pour (a), $\sigma u \sigma^{-1}$ et $u^{-1}$ envoient tous deux $s_0$ sur $\lambda^{-1}s_0$ et $t_0$ sur $\lambda t_0$, donc sont égaux. Pour (b), si $u \in U_0$ et $t$ correspond à $\lambda$, $tut^{-1}s_0 = t\lambda^{-1}us_0 = \lambda^{-1}ts_0 = s_0$. Pour (c), soit $a t b = a' t' b'$. On pose $c' = t'^{-1}a'^{-1}at' \in U_0$, $r = t'^{-1}t \in T_H$ et $d = bb'^{-1} \in U_1$ ; alors $c'rd = 1$. Si $r$ correspond à $\lambda$, $c'^{-1} = rd$ fixe $s_0$ et envoie $t_0$ sur $rt_0 = \lambda^{-1}t_0$ ; la fidélité donne $\lambda = 1$, donc $r$ fixe $s_0$ et $t_0$ et $r = 1$. Il reste $c'd = 1$, avec $d \in U_0 \cap U_1$, d'où $d = 1$, $c' = 1$, et $(a, t, b) = (a', t', b')$.

Point~2. L'ensemble $Z_T$ est un sous-groupe : si $\lambda s_0 = ts_0$ et $\mu s_0 = t's_0$, alors $\lambda\mu s_0 = t'ts_0$ et $\lambda^{-1}s_0 = t^{-1}s_0$. De (a) on tire $\sigma t^{-1} = t\sigma$, d'où $\lambda^{-1}t_0 = \sigma\lambda^{-1}s_0 = \sigma t^{-1}s_0 = t t_0$ : l'arc $(\lambda hs_0, \lambda^{-1}ht_0)$ est l'image par $ht$ de $(s_0, t_0)$. Pour la fidélité, soit $u \in U_0$ avec $ut_0 = \lambda t_0 = \sigma t s_0$. Alors $a = t^{-1}\sigma^{-1}u\sigma \in U_0$, et $b = \sigma^{-1}u\sigma$ est dans $U_1$ puisque $\sigma^{-1}b\sigma = u$. L'égalité $(tat^{-1})\, t \cdot 1 = 1 \cdot 1 \cdot b$ et l'injectivité (c) donnent $t = 1$, donc $\lambda s_0 = s_0$ et $\lambda = 1$.

Point~3. Si $\lambda \in Z$, un $u$ associé à $\lambda$ par la condition d'homothétie est dans $T_H$ et $\lambda s_0 = us_0$, donc $\lambda \in Z_{T_H}$ ; réciproquement, si $\lambda \in Z_{T_H}$, $\lambda s_0 = ts_0 = \mu s_0$ pour un $\mu \in Z$, et $\lambda = \mu$. Si $t \in T$, la permutation $hU_0 \mapsto htU_0$, bien définie parce que $t$ et $t^{-1}$ normalisent $U_0$, est dans $Z_T$ et lui associe $t$, donc $t \in T_{Z_T}$ ; réciproquement, si $u$ est associé à $\lambda \in Z_T$ avec $\lambda s_0 = ts_0$, $u$ et $t$ coïncident sur $s_0$ et sur $t_0$, donc $u = t$.
\end{proof}

\begin{remarque}
Les deux sens de l'équivalence tiennent, y compris la réciproque que la page~247 vérifie sur des lignes en partie illisibles. La commutativité de $Z$, que la piste suppose et que la page met entre crochets, découle de la condition d'homothétie et de la fidélité (72), comme le dit la marge \q{[ ] résulte des autres} ; la démonstration n'utilise pas la simple transitivité. La fidélité de $H$ sur $H/U_0$, condition (77), n'intervient nulle part, ni $s_0 \neq t_0$. La condition d'homothétie est prise sous la forme corrigée par la lecture, $u(s_0, t_0) = (\lambda s_0, \lambda^{-1}t_0)$. L'isomorphisme $T_H \simeq Z$ de (81') n'est pas formalisé comme morphisme de groupes, seulement la bijection entre les deux familles.
\end{remarque}

\begin{encadre}
\begin{theoreme}[Grothendieck, pages 92 à 95, le commutant\piste{79-commutant-octahedral-exception}]
Soit $R$ un anneau commutatif.
\begin{enumerate}
\item Si $R$ est exceptionnel, l'application $\sigma(\dot x) = (x + 2y)^{\cdot}$, pour un $y$ tel que $x \wedge y = 1$, est bien définie sur $S_M$. Elle commute à $\mathrm{SL}(2, R)$ et aux scalaires, et $\sigma^2 = 1$. Si de plus $2 \neq 0$, $\sigma(v) \neq cv$ pour tout $c \in R^*$ et tout $v$ ; en particulier $\sigma$ n'a pas de point fixe.
\item Les scalaires commutent à $\mathrm{SL}(2, R)$, et $c$, $c'$ agissent de même si et seulement si $c' = \pm c$.
\item Toute application $f : S_M \to S_M$ qui commute à $\mathrm{SL}(2, R)$ est un scalaire, ou bien $R$ est exceptionnel et $f = c\,\sigma$ pour un $c \in R^*$.
\item Sur $\mathbf Z/n$, $2R = \{0, 2\} \neq \{0\}$ si et seulement si $n = 4$ ; pour $n \neq 4$ le commutant est formé des scalaires.
\item Sur $\mathbf Z/4$, si $x \wedge y = 1$, les $z \in M^*$ tels que $x \wedge z$ ne soit pas inversible sont $\pm x$ et $\pm(x + 2y)$ : $\sigma(\dot x)$ est le seul sommet autre que $\dot x$ qui ne lui soit pas adjacent, l'antipode sur l'octaèdre.
\end{enumerate}
\end{theoreme}
\end{encadre}

\begin{proof}
Si $x \wedge y = 1$, tout $z$ s'écrit $z = (z \wedge y)\,x + (x \wedge z)\,y$. Dans un anneau exceptionnel, $4 = 0$ (car $2 \cdot 2 \in \{0, 2\}$) et $2c = 2$ pour tout $c$ inversible.

Point~1. Deux partenaires de $x$ diffèrent d'un multiple $tx$ ; comme $2t \in \{0, 2\}$ et $4 = 0$, $x + 2y' = x + 2y$ ou $x + 2y' = -(x + 2y)$. Le passage de $x$ à $-x$ change $y$ en un partenaire de $-x$ et donne de même $\pm(x + 2y)$ ; $\sigma$ est donc bien définie. Si $g \in \mathrm{SL}(2, R)$, $gy$ est un partenaire de $gx$, d'où $\sigma g = g\sigma$. Comme $(x + 2y) \wedge y = 1$, $\sigma^2(\dot x) = (x + 4y)^{\cdot} = \dot x$. Pour $c \in R^*$, $c^{-1}y$ est un partenaire de $cx$ et $2c^{-1} = 2c = 2$, d'où $\sigma(c\dot x) = c\,\sigma(\dot x)$. Si $x + 2y = \pm cx$, on applique $x \wedge \cdot$ : $2 = 0$.

Point~3. Soit $v = (p, q)$ un représentant de $f(\dot e_1)$, avec $e_1 = (1, 0)$. Les matrices de lignes $(1, t)$ et $(0, 1)$ fixent $\dot e_1$, donc $(p + tq, q) = \pm(p, q)$ pour tout $t$. Pour $t = 1$ : ou bien $q = 0$, et $p$ est inversible puisque $v \in M^*$ ; ou bien $2q = 0$ et $q = -2p$. Dans ce second cas $p(s + 2r) = 1$ si $v \wedge (r, s) = 1$, donc $p$ est inversible. Pour $a \in R$ et $t = ap^{-1}$, $tq = -2a$, et la condition donne $2a = 0$ ou $2a = 2p$. Avec $a = 1$, $2p = 2$ ; donc $2a \in \{0, 2\}$ pour tout $a$, $R$ est exceptionnel, et $q = -2p = 2$. Pour $\dot x$ quelconque, soit $g$ la matrice de colonnes $x$ et $y$, avec $x \wedge y = 1$ ; alors $g\dot e_1 = \dot x$ et $f(\dot x) = g f(\dot e_1) = (px + qy)^{\cdot}$. Si $q = 0$, c'est $p\dot x$. Sinon c'est $(px + 2y)^{\cdot} = (p(x + 2y))^{\cdot} = p\,\sigma(\dot x)$, puisque $2p = 2$.

Point~4. Si $2\cdot 2 \in \{0, 2\}$ dans $\mathbf Z/n$, $n$ divise $4$ ou $2$ ; on écarte $n = 1, 2$ par $2 \neq 0$ et l'on vérifie $n = 4$ par calcul. Si $R$ n'est pas exceptionnel, le point~3 donne un scalaire ; si $2 = 0$, $\sigma$ est l'identité. Point~5. On écrit $z = ax + by$ avec $b = x \wedge z$ non inversible, donc $b \in \{0, 2\}$. Comme $z \in M^*$, $a$ est inversible, donc $a = \pm 1$, et $-x + 2y = -(x + 2y)$.
\end{proof}

\begin{remarque}
L'énoncé tient sur tout anneau commutatif, sans l'hypothèse $2 \neq 0$ : quand $2 = 0$, $R$ est exceptionnel au sens adopté ici mais $\sigma$ est l'identité, et le commutant se réduit aux scalaires. L'application $f$ n'a pas à être supposée bijective. Le produit direct de la page est confirmé : $\sigma$ commute aux scalaires, $\sigma^2 = 1$, et $\sigma \neq c$ pour tout $c$, de sorte que le commutant est $R^*/\pm 1 \times \{1, \sigma\}$. La démonstration ne passe pas par le normalisateur $N_{G'}(H')$ : elle fixe l'image de $\dot e_1$, et le calcul est le même. La formule biffée de la page~95, $\sigma(\dot x) = (x + 2y)^{\cdot}$, est celle que l'on démontre ; l'élément $\sigma$ envoie $\dot e_1$ sur la classe de $(1, 2)$, image de $e_1$ par la matrice de lignes $(1, 0)$ et $(2, 1)$.
\end{remarque}

\begin{enonce}[Pages 275 à 279, un sommet et son étoile\piste{79-vertex-determined-by-star}]
Soient $R$ un anneau commutatif, $x, y \in M$ avec $x \wedge y = 1$, et $x' \in M$.
\begin{enumerate}
\item On a l'équivalence : pour tout $z$ tel que $x \wedge z$ soit inversible, $x' \wedge z = \pm\, x \wedge z$, si et seulement si $x' = \pm x$, ou bien $R$ est exceptionnel et $x' = \pm(x + 2y)$, c'est-à-dire $\dot x' = \sigma(\dot x)$.
\item Si $R$ est exceptionnel, la matrice $A$ de lignes $(1, 2)$ et $(2, -1)$ vérifie $Ax = \pm(x + 2y)$ ; elle a pour déterminant $-1$ et $A^2 = 1$.
\end{enumerate}
\end{enonce}

\begin{proof}
On écrit $x' = \alpha x + \beta y$, avec $\alpha = x' \wedge y$ et $\beta = x \wedge x'$. Pour $s \in R$, $z = y + sx$ vérifie $x \wedge z = 1$ et $x' \wedge z = \alpha - s\beta$. L'hypothèse donne $\alpha - s\beta = \pm 1$ pour tout $s$. Avec $s = 0$, $\alpha = \pm 1$. Si $\alpha = 1$, $s = 1$ donne $\beta = 0$, d'où $x' = x$, ou $\beta = 2$ ; dans ce dernier cas, $1 - 2s = \pm 1$ donne $2s \in \{0, 2\}$ pour tout $s$, et $x' = x + 2y$. Le cas $\alpha = -1$ donne de même $x' = -x$, ou $\beta = -2$ et $x' = -(x + 2y)$. Réciproquement, $(x + 2y) \wedge z = x \wedge z - 2(z \wedge y)$ ; si $2(z \wedge y) = 0$ c'est $x \wedge z$, et si $2(z \wedge y) = 2$ c'est $x \wedge z - 2 = -x \wedge z$, puisque $2(x \wedge z) = 2$ pour $x \wedge z$ inversible.

Pour le point~2, on note $x = (p, q)$, $y = (r, s)$, avec $ps - qr = 1$. Chacun de $2p, 2q, 2r, 2s$ vaut $0$ ou $2$, et $2(ps) = (2p)s$ vaut $0$ si $2p = 0$, $2s$ si $2p = 2$. On examine les seize cas ; l'égalité $2(ps) - 2(qr) = 2$ en écarte une partie, et dans les autres on a $2q = 2r$ et $2p = 2q + 2s$, ce qui donne $Ax = x + 2y$, ou bien $2p + 2q + 2r = 0$ et $2p = 2s$, ce qui donne $Ax = -(x + 2y)$. Enfin $\det A = -1 - 4 = -1$ et $A^2 = 5 \cdot 1 = 1$.
\end{proof}

\begin{remarque}
L'hypothèse \q{semi-local} de la page n'est pas nécessaire : l'énoncé tient sur tout anneau commutatif, et la démonstration n'utilise ni les corps résiduels ni le lemme de colinéarité de la page~276. Il suffit des voisins $y + sx$ d'un seul type, ceux pour lesquels $x \wedge z = 1$. L'énoncé dit aussi que, dans le cas exceptionnel, le seul sommet autre que $\dot x$ qui ait la même étoile est $\sigma(\dot x)$ : la question que la page laisse ouverte (\q{il n'est pas vrai nécessairement \dots}) est donc tranchée dans ce cadre algébrique, où l'étoile est prise type par type. La transcription signale à la page~277 la conclusion $\alpha - \beta \in \pm(\alpha - \beta)$, qui doit se lire $\alpha + \beta \in \pm(\alpha - \beta)$ ; la démonstration ci-dessus ne passe pas par elle. L'involution $a$ de la page~278 est l'élément $\sigma$ de l'énoncé précédent, et la matrice $A$ réalise le déterminant $-1$ annoncé, sur tout anneau exceptionnel ; pour $\mathbf Z/4$ c'est la matrice que la lecture vérifie. Le lemme combinatoire de la page~275 (deux sommets distincts de même étoile forment un double cône) n'est pas formalisé.
\end{remarque}

\begin{enonce}[Pages 19 et 20, distance au plus $3$\piste{79-semilocal-diameter-three}]
Soit $R$ un anneau commutatif qui n'a qu'un nombre fini d'idéaux maximaux.
\begin{enumerate}
\item Si $a_0d + b_0d' = 1$, il existe $a \in R$ tel que $ad + d'$ soit inversible.
\item Pour tout $e \in R^*$ et tous $x, y \in M^*$, il existe $z_1, z_2 \in M^*$ avec $x \wedge z_1 = z_1 \wedge z_2 = z_2 \wedge y = e$ : deux sommets de $SP_M$ sont à distance au plus~$3$.
\end{enumerate}
Cela vaut en particulier pour un anneau local, un corps, et $\mathbf Z/n$ pour $n \geq 1$.
\end{enonce}

Le point~2 donne un chemin de trois arêtes de type $e$ entre deux sommets quelconques :
\[
\begin{tikzcd}[column sep=large]
\dot x \arrow[r, dash, "e"] & \dot z_1 \arrow[r, dash, "e"] & \dot z_2 \arrow[r, dash, "e"] & \dot y
\end{tikzcd}
\qquad x \wedge z_1 = z_1 \wedge z_2 = z_2 \wedge y = e.
\]

\begin{proof}
Point~1. L'égalité $a_0d + b_0d' = 1$ dit que $d'$ et $d$ sont étrangers ; par le rang stable~$1$ des anneaux semi-locaux, énoncé général rappelé à la suite de la démonstration, $d' + ad$ est inversible pour un $a \in R$.

Point~2. Soient $w$ et $y'$ avec $x \wedge w = 1$ et $y \wedge y' = 1$. L'identité $(x \wedge y)(w \wedge y') - (x \wedge y')(w \wedge y) = (x \wedge w)(y \wedge y')$ donne $(w \wedge y')\,d - (x \wedge y')\,d' = 1$ pour $d = x \wedge y$ et $d' = w \wedge y$. Par le point~1, $u = ad + d'$ est inversible pour un $a$. Le vecteur $z = ax + w$ vérifie $x \wedge z = 1$ et $z \wedge y = u$. On pose $z_1 = ez$, d'où $x \wedge z_1 = e$ et $z_1 \wedge y = eu$, puis, suivant la page~19, $z_2 = u^{-1}(z_1 + y)$ : $z_1 \wedge z_2 = u^{-1}eu = e$ et $z_2 \wedge y = e$. Les vecteurs $z_1$ et $z_2$ sont dans $M^*$, puisque $z_1 \wedge y$ et $z_2 \wedge y$ sont inversibles. Un anneau local n'a qu'un idéal maximal ; un anneau fini n'a qu'un nombre fini d'idéaux.
\end{proof}
\general{SemilocalStableRange}

\begin{remarque}
L'énoncé tient sur tout anneau semi-local, au sens où l'ensemble des idéaux maximaux est fini. Pour $\Lambda = \mathbf Z/n$ et $n \geq 5$, le diamètre~$3$ du graphe est le théorème~11 de Singerman et Strudwick, mis en ligne en 2018 et publié en 2020 \cite{singermanstrudwick2020} (numérotation de la version arXiv). La page cite comme connu (\q{on sait que}) le rang stable~$1$ des anneaux semi-locaux, dû à Bass (1964, §~4 et lemme~6.4) \cite{bass1964} ; le fichier général le démontre par le théorème chinois. La démonstration ne construit pas un partenaire commun dans chaque corps résiduel comme le propose la lecture : elle cherche ce partenaire sur la droite $ax + w$, où $x \wedge z = 1$ est automatique. La borne est atteinte sur $\mathbf Z/n$ dès que $(\mathbf Z/n)^*$ contient un $\mu \neq \pm 1$, c'est-à-dire pour $n \notin \{1, 2, 3, 4, 6\}$, en particulier pour l'icosaèdre ($n = 5$) : c'est le point~4 de l'énoncé de la page~188 ci-dessous. Le cas $n = 6$, que couvre le théorème de Singerman et Strudwick pour $n \geq 5$, n'est pas formalisé. Le cas $\Lambda = \mathbf Z$, qui n'est pas semi-local, n'est pas traité.
\end{remarque}

\begin{enonce}[Pages 96 à 101 et 136 à 140, la torsion dans le modèle\piste{79-star-polygon-twist}]
Soient $R$ un anneau commutatif, $e, \mu \in R^*$ et $x, y \in M$ avec $x \wedge y = e$.
\begin{enumerate}
\item (Torsion d'une extrémité.) $x \wedge \mu y = \pm \mu e$ si et seulement si $x \wedge y = \pm e$ : les arêtes $\{\dot x, \mu\dot y\}$, pour $\{\dot x, \dot y\}$ arête de type $\pm e$, sont les arêtes de type $\pm\mu e$. De même $\mu x \wedge z = \pm\mu e$ si et seulement si $x \wedge z = \pm e$ : l'étoile de $\dot x$ dans $P_{M, \pm e}$ est, comme partie de $S_M$, l'étoile de $\mu\dot x$ dans $P_{M, \pm\mu e}$.
\item (L'étoile est une droite.) Si $x \wedge z = e$, alors $z = y + kx$ avec $k = z \wedge e^{-1}y$.
\item (Formule (178), page~67, \lot{79}{4}{67}.) Si $x \wedge z = \pm e$ et $z \wedge y = \pm e$, alors $z = ax + by$ avec $a, b \in \{1, -1\}$.
\item (Le polygone étoilé.) Pour $k, k' \in R$, $(y + kx) \wedge (y + k'x) = (k - k')\,e$ ; ces deux points sont joints par une arête de type $\pm\mu e$ si et seulement si $k - k' = \pm\mu$.
\item (L'indice.) $\mu(y + kx) = \mu y + (\mu k)x$, et $x \wedge \mu y = \mu\,(x \wedge y)$.
\end{enumerate}
\end{enonce}

\begin{proof}
Les points~1 et~5 viennent de la bilinéarité : $x \wedge \mu y = \mu\, x \wedge y$, et l'on simplifie par $\mu$ inversible. Pour le point~2, $(x, e^{-1}y)$ est une base de déterminant~$1$ et $z = (z \wedge e^{-1}y)\,x + (x \wedge z)\,e^{-1}y$. Le point~3 se lit sur la même décomposition : les coefficients sont $e^{-1}(z \wedge y)$ et $(x \wedge z)e^{-1}$, égaux à $\pm 1$. Pour le point~4, $(y + kx) \wedge (y + k'x) = k'(y \wedge x) + k(x \wedge y) = (k - k')e$, et $(k - k')e = \pm\mu e$ équivaut à $k - k' = \pm\mu$.
\end{proof}

\begin{remarque}
Ces identités vérifient le calcul de la lecture dans le modèle, et non l'énoncé abstrait de la piste. Sur l'étoile de $\dot x$, indexée par $k \mapsto (y + kx)^{\cdot}$, le polyèdre $P_{M, \pm e}$ induit le polygone $k \sim k \pm 1$ et $P_{M, \pm\mu e}$ le polygone $k \sim k \pm \mu$, sa puissance $\mu$-ième. La multiplication par $\mu$ envoie l'indice $k$ sur $\mu k$ : l'exposant est $\mu$ lui-même, sans la puissance illisible des formules (65) et (66), comme la lecture l'avait calculé. Les identités tiennent sur tout anneau commutatif, et pas seulement sur $\mathbf Z/n$. Ne sont pas formalisés : l'action fidèle de $(\mathbf Z/n)^*/\pm 1$ sur les structures de polygone d'un ensemble quelconque à $n$ éléments, l'axiome To~6 dans le cadre abstrait et l'injection $Z \hookrightarrow (\mathbf Z/n)^*/\pm 1$ qui en découle hors de l'octaèdre, et la comparaison avec les opérateurs de Wilson.
\end{remarque}

\begin{encadre}
\begin{theoreme}[Page 188, forme de l'édition, distance à l'image\piste{79-homothety-distance-three}]
\begin{enumerate}
\item (Cadre abstrait.) Soient $H$, $U_0$, $\sigma$ comme ci-dessus, avec $\sigma^2 \in U_0$, et $\lambda \neq 1$ une permutation de $S = H/U_0$ qui commute à $H$. Alors $\lambda$ n'a pas de point fixe ; si $\lambda s$ est voisin de $s$, c'est le seul voisin de $s$ ; si $s$ et $\lambda s$ ont un voisin commun, ils ont les mêmes voisins. En particulier, si $s$ a au moins deux voisins et si aucun autre sommet n'a les voisins de $s$, alors $\lambda s$ n'est ni $s$, ni un voisin de $s$, ni à distance~$2$ de $s$ ; s'il est à distance au plus~$3$, il est à distance~$3$.
\item (Modèle, sur un anneau commutatif $R$ non nul.) Soient $e \in R^*$ et, sur $S_M$, le graphe des arêtes de type $\pm e$ : $\dot x$ et $\dot z$ sont adjacents si $x \wedge z = \pm e$. Soit $f \neq \mathrm{id}$ une application de $S_M$ dans lui-même qui commute à $\mathrm{SL}(2, R)$. Pour tout sommet $v$, $f(v) \neq v$, $f(v)$ n'est pas adjacent à $v$, et $v$ et $f(v)$ ont un voisin commun si et seulement si $R$ est exceptionnel et $f = \sigma$. Pour $f = \sigma$, un voisin commun de $\dot x$ et $\sigma(\dot x)$ est $(ey)^{\cdot}$, si $x \wedge y = 1$.
\item Si de plus $R$ n'a qu'un nombre fini d'idéaux maximaux et si $f \neq \sigma$, chaque sommet $v$ est à distance~$3$ de $f(v)$.
\item Sur $\mathbf Z/n$, $n \geq 2$, pour $\mu \in (\mathbf Z/n)^*$ avec $\mu \neq \pm 1$ et tout type $\pm e$, chaque sommet $\dot x$ est à distance~$3$ de $\mu\dot x$.
\end{enumerate}
\end{theoreme}
\end{encadre}

\begin{proof}
Point~1. Une permutation qui commute à $H$ et fixe un sommet $hs_0$ fixe $s_0$, donc est l'identité. Les arcs sont permutés par $H$, et $\sigma^2 \in U_0$ rend le graphe non orienté : $h\sigma$ envoie $(s_0, t_0)$ sur $(ht_0, hs_0)$. Si $(s, t)$ et $(s, t')$ sont des arcs, images de $(s_0, t_0)$ par $k$ et $k'$, alors $g = k'k^{-1}$ fixe $s$ et envoie $t$ sur $t'$ : le stabilisateur de $s$ est transitif sur les voisins de $s$. Il fixe aussi $\lambda s$, puisque $\lambda(gs) = g\lambda s$. Si $\lambda s$ est voisin de $s$ et $t'$ un autre voisin, on prend $g$ fixant $s$ avec $g\lambda s = t'$, d'où $t' = \lambda s$. Si $t$ est voisin de $s$ et de $\lambda s$, et $t'$ voisin de $s$, on prend $g$ fixant $s$ avec $gt = t'$ ; l'arête $\{t, \lambda s\}$ a pour image $\{t', \lambda s\}$. Dans l'autre sens, si $t'$ est voisin de $\lambda s$, on prend $g$ fixant $\lambda s$ avec $gt = t'$ ; alors $\lambda(gs) = \lambda s$, donc $gs = s$ par injectivité de $\lambda$, et l'arête $\{s, t\}$ a pour image $\{s, t'\}$. Les deux hypothèses du point~1 excluent alors les distances $0$, $1$ et~$2$.

Point~2. Par l'énoncé du commutant, $f$ est un scalaire $c$, ou bien $R$ est exceptionnel et $f = c\sigma$. Si $f = c$, $f \neq \mathrm{id}$ donne $c \neq \pm 1$. Si $cx = \pm x$, on applique $\cdot \wedge y$ : $c = \pm 1$, exclu. On a $x \wedge cx = 0 \neq \pm e$, puisque $R$ est non nul. Si $x \wedge z = \varepsilon_1 e$ et $z \wedge cx = \varepsilon_2 e$, avec $\varepsilon_i = \pm 1$, alors $-c\varepsilon_1 e = \varepsilon_2 e$, et $c = \pm 1$ en simplifiant par $e$. Si $f = c\sigma$ et $2 = 0$, $\sigma$ est l'identité et l'on est dans le cas précédent. Supposons $2 \neq 0$, de sorte que $f(\dot x) = (c(x + 2y))^{\cdot}$ avec $2c = 2$. Si $c(x + 2y) = \pm x$, on applique $x \wedge \cdot$ : $2 = 0$, exclu. On a $x \wedge c(x + 2y) = 2c = 2$ ; si $2 = \pm e$, $2$ est inversible, et $2 \cdot 2^{-1} \in \{0, 2\}$ donne $1 = 0$ ou $1 = 2$, impossible dans un anneau non nul. Soit enfin $z$ voisin de $\dot x$ et de $f(\dot x)$, avec $x \wedge z = \varepsilon_1 e$ et $b = z \wedge y$. Alors $z \wedge c(x + 2y) = c(-\varepsilon_1 e + 2b)$, avec $2b \in \{0, 2\}$ et $2 = 2e$ ; comme $4 = 0$, $-\varepsilon_1 e + 2b = \pm e$ dans tous les cas, et l'égalité $c(\pm e) = \varepsilon_2 e$ donne $c = \pm 1$, donc $f = \sigma$. Réciproquement, $x \wedge ey = e$ et $ey \wedge (x + 2y) = -e$.

Points~3 et~4. Sur un anneau semi-local, deux sommets quelconques sont joints par un chemin de longueur~$3$ formé d'arêtes de type $e$ (énoncé des pages~19 et~20). Sur $\mathbf Z/n$, $n \geq 2$, l'anneau est fini et non nul, et l'on applique le point~2 aux scalaires.
\end{proof}

\begin{remarque}
La partie abstraite de la piste est démontrée dans le cadre du graphe de classes, sous la forme : pas de point fixe, pas de voisin, et un voisin commun force l'égalité des étoiles. Elle n'utilise que la transitivité de $H$ sur les arcs et la commutation de $\lambda$ à $H$ ; ni la simple transitivité, ni la condition d'homothétie, ni la structure de surface n'interviennent. L'exclusion de la distance~$1$, que la page attribue aux stabilisateurs (lecture incertaine), demande que $s$ ait au moins deux voisins. Le passage de \q{même étoile} à \q{octaèdre} est le lemme du double cône de la page~275, qui n'est pas formalisé ; l'énoncé abstrait le remplace par l'hypothèse qu'aucun autre sommet n'a l'étoile de $s$. La page parle de $\lambda^{-1}s$ ; comme $\lambda^{-1}$ commute aussi à $H$, cela revient au même, et l'énoncé vaut pour un sommet comme pour tous (\q{c'est kif-kif}).

L'énoncé de la page est celui de l'encadré non vérifié ci-dessous ; le cadre abstrait du point~1 et l'application au modèle des points~2 à~4 sont de l'édition. Dans le modèle, l'énoncé tient sur tout anneau commutatif non nul pour les distances $0$, $1$ et~$2$, pour toute application qui commute à $\mathrm{SL}(2, R)$ et pour tout type $\pm e$ ; pour les scalaires le calcul est direct ($x \wedge \mu x = 0$, et $\mu e = \pm e$ force $\mu = \pm 1$), sans argument de stabilisateur. L'exception de la distance~$2$ est $\sigma$, et elle seule. Sur $\mathbf Z/n$ elle ne se produit que pour $n = 4$, où $\sigma$ est l'antipodie de l'octaèdre et où les scalaires sont triviaux. Sur un autre anneau exceptionnel avec $2 \neq 0$, par exemple $\mathbf Z/4 \times \mathbf F_2$, $\sigma$ est aussi à distance~$2$ ; un calcul hors de Lean, à partir du polygone de l'étoile $k \sim k \pm 1$ indexé par $R$, montre que l'étoile y est formée de deux cycles de longueur~$4$ ; le polyèdre n'est donc pas une surface, et l'exemple ne contredit pas la page. Pour $e = 1$ et $n \geq 5$, le point~4 est le théorème~9 de Singerman et Strudwick sur les pôles de la carte de Farey modulo~$n$ \cite{singermanstrudwick2020}.
\end{remarque}

Ne sont pas démontrés ici : l'isomorphisme de groupes $T_H \simeq Z$ de (81') et le cas $\mathfrak z_H \neq 1$ de (88) ; le lemme combinatoire du double cône (page~275), qui ramènerait l'hypothèse \q{aucun autre sommet n'a la même étoile} de l'énoncé de la page~188 à l'exclusion de l'octaèdre dans un 2-polyèdre régulier abstrait ; l'optimalité de la borne~$3$ pour $\mathbf Z/6$ et pour un anneau semi-local quelconque ; la connexité de $SP_M$ pour $\Lambda = \mathbf Z$ ; l'action de $(\mathbf Z/n)^*/\pm 1$ sur les structures de polygone, l'axiome To~6 abstrait et l'injection $Z \hookrightarrow Z'_n$ (pages~136 à~140) ; la question de la page~141 sur les triangulations dont les sommets sont d'ordre~$7$ ; la distance exactement~$3$ de la page~188 hors des anneaux semi-locaux.

\begin{encadre}
\begin{theoreme}[Grothendieck, pages 187 et 188, homothéties à distance trois\piste{79-homothety-distance-three}]\nonverifie
Dans un 2-polyèdre régulier triangulé abstrait, deux sommets distincts n'ont la même étoile que dans l'octaèdre (lemme du double cône, page~275). Une homothétie $\lambda \neq \mathrm{id}$ n'envoie donc aucun sommet à distance $1$ ou $2$, l'antipodie de l'octaèdre exceptée, et quand $\lambda^{-1}s$ est à distance au plus $3$ de $s$, cette distance vaut exactement~$3$.
\end{theoreme}
\end{encadre}

\begin{encadre}
\begin{theoreme}[Grothendieck, pages 136 à 140 et 184 à 187, polygones étoilés\piste{79-star-polygon-twist}]\nonverifie
Le groupe $(\mathbf Z/n)^*/\pm 1$ opère sur les structures de polygone d'un ensemble à $n$ éléments. Sous l'axiome To~6, une homothétie agit sur une triangulation en redessinant chaque étoile comme le polygone étoilé d'indice $\nu \in (\mathbf Z/n)^*/\pm 1$, recentré au sommet image, avec $\nu$ indépendant du sommet, et l'on a une injection $Z \hookrightarrow Z'_n$, sauf pour l'octaèdre.
\end{theoreme}
\end{encadre}
